% =====================================================================
%  §6 — Discussion
%
%  Revised 2026-04-21: subsections collapsed from four to two.
%  §6.1 (positioning) + §6.2 Model Extensions → §6.1 Discussion.
%  §6.3 Limitations + §6.4 Future Work → §6.2 Limitations and
%  Future Work. Labels preserved for all cross-section references
%  (sec:disc-future is referenced from §2).
% =====================================================================

\section{Discussion}
\label{sec:discussion}

Two features of the proposal merit brief emphasis. First, the
always-required vocabulary is small, and conditional activation keeps
the extended vocabulary compact without sacrificing expressive power.
Second, every assertion is anchored to a specific source span in its
source publication, making each annotation independently auditable.

\subsection{Positioning and Model Extensions}
\label{sec:disc-positioning}
\label{sec:disc-extensions}

The released schema formalises the implemented conformance
requirements as SHACL shapes, makes the implemented activation
conditions machine-checkable, and anchors each assertion to its
source text rather than only to a style guide. The flag vocabularies
of \cref{sec:workflow-flags} make provenance explicit: a reader can
tell whether a value was drawn from the author's summary, proposed
by a reviewer, added as an AI assumption, or left deliberately
unpopulated. More broadly, the schema is intended to make literature review tractable as
a \emph{semantic-parsing} problem: the mapping of natural-language claims
to a shared, source-anchored representation. This is not a task a language model can perform in isolation. The target of
each parse is a structured, source-anchored representation, so the quality
of the parse depends on the quality of that target language. We see the present
contribution as one increment along a longer path, building on
decades of biomedical-ontology work (ECO, SEPIO, HPO, IAO, OBI, VRS,
and others), not as a replacement for any of it. A complementary
downstream layer, in which an epistemological type system and
machine-checkable meta-predicates constrain the decision logic that
consumes such evidence, is addressed in~\cite{our_preprint}.

The pilot identified two candidate dimensions that were exercised by
multiple papers and promoted to first-class status, Phenotype Scale and
Variant Ascertainment. It also identified fifteen further candidate
extensions that remain unpromoted under the two-papers rule
(\cref{sec:workflow-promotion}). The pattern of
candidates is summarised in \cref{tab:candidates}; per-candidate
rationale, affected dimensions, and proposed enumeration values
are documented in a companion \texttt{schema/EXTENSIONS.md} in the
repository.\footnote{\url{https://github.com/ForomePlatform/genetic-evidence-model/blob/master/schema/EXTENSIONS.md}}
The Gupta annotation, which identified no candidate extensions, is also a
positive finding. It reports a family with an X-linked splice-site mutation
and careful pedigree segregation, a canonical evidence shape for which the
model was designed. The schema accommodates this evidence without strain.

\begin{table}[tbp]
\centering
\caption{Unpromoted candidate schema extensions surfaced by the
six-paper pilot, grouped by concern. Per-candidate detail is in the
companion \texttt{schema/EXTENSIONS.md}; paper keys are J~=~Jossin,
D~=~Davis, N~=~Nelson, DU~=~Duerr, IN~=~Inouye. Gupta surfaced none.}
\label{tab:candidates}
\small
\begin{tabularx}{\linewidth}{@{}llX@{}}
\toprule
Group & Papers & Count and examples \\
\midrule
Finer structural resolution  & J, N   & 4 (subdomain, interaction, gene-relation, cross-level variant) \\
Polygenic-score machinery    & IN     & 6 (score target, whole-genome aggregate, target composition, epidemiological measurement, cohort sub-block, derived-artefact flag) \\
Effect-size and curation     & DU, D  & 3 (direction of effect, knowledge-domain priority, curator critique) \\
Cross-item and activation scope & DU & 2 (cross-item/translational synthesis, conditional-activation scope) \\
\bottomrule
\end{tabularx}
\end{table}

\subsection{Limitations and Future Work}
\label{sec:disc-limitations}
\label{sec:disc-future}

The pilot has significant limitations. Six papers and a single
annotator do not support claims about inter-annotator agreement or
generalisability; we treat the work as a feasibility study, not a
benchmark. The evaluation is not a clinical
deployment: although the model is motivated by variant
interpretation in clinical-genomics programs, the pilot stops short
of applying the model to a live workflow.

Three near-term directions follow from these limitations. First, annotate a
second corpus selected to exercise the strongest candidate extensions.
Second, generate OWL, SHACL, JSON Schema, and LinkML renderings from a
single source of truth. This would open the schema to the OBO and Bridge2AI
ecosystems and support term-level alignment of the dimension vocabularies
with relevant ontologies, for example by aligning the molecular and cellular
phenotype scales with Gene Ontology biological processes. Third, conduct a
clinical-deployment study to test whether the feasibility claim translates
into utility.

A further limitation revealed by the UMLS crosswalk is representational
rather than empirical. Several GEM tokens are intersective compounds:
\dimval{statistical\_genetics} denotes statistics~$\times$~genetics
(\texttt{C0038215}~$\times$~\texttt{C0017398}), and
\dimval{bioinformatics\_inference} denotes inference~$\times$~bioinformatics.
Relational tokens have the same structure. For example,
\dimval{related\_gene} combines a relation operator with \emph{Genes}
(\texttt{C0017337}). Because the current crosswalk maps each GEM token to
one UMLS concept, it must sometimes use the nearest available anchor rather
than a concept that expresses the full compound meaning. For
\dimval{statistical\_genetics}, that anchor is \emph{Biostatistics},
recorded with the relation \emph{narrower}.

Concept-level post-coordination would address this limitation by mapping a
token to a small set of concepts whose intersection represents its intended
meaning. The UMLS Metathesaurus does not provide such a mechanism at the
aggregate level. As an aggregator, it records the concepts asserted by its
source vocabularies and the links among them, but it does not define a
compositional grammar of its own. Composition remains a feature of individual
source vocabularies, with the compositional grammar of SNOMED~CT being the
best-known example. It cannot be transferred directly across the aggregate:
the same intended intersection may be pre-coordinated in one source, absent
from another, and decomposed differently in a third~\cite{umls2004}.

The problem is particularly clear for \dimval{related\_gene}. UMLS often
contains pre-coordinated specific instances while lacking a generic concept:
gene-relatedness appears in terms such as \emph{modifier},
\emph{homologous genes}, and \emph{linked genes}, whereas the generic sense
of ``related in some way'' is represented only by the bare operator
\emph{Associated with} (\texttt{C0332281}). Thus
\dimval{related\_gene} can be represented as
\texttt{C0332281}~$\times$~\texttt{C0017337}, but not by a single concept.
Extending the crosswalk with an explicit composite relation, represented as
a structured set of concepts, is future work. At present, the intended composition
is preserved in the mapping's curator note; where the relation is the
semantic core, the operator concept itself is recorded as the anchor with
relation \emph{related}.

The crosswalk review also highlighted a representational limitation of the
Resolution axis. Its current enumeration presupposes genomic coordinates, but
genetic evidence may be localized in several coordinate systems. A variant
may be described at the genomic level (HGVS~\texttt{g.}), the transcript
level (\texttt{c.}), or the protein level (\texttt{p.}).

For example, an assertion of the form ``variant $V$ has effect $E$ in
transcript $T$'' is resolved in transcript coordinates. Such assertions are
often tissue-restricted and are characteristic of bioinformatics inference.
The domain mapping underlying the \dimval{PROTEIN\_SUBDOMAIN} candidate is
instead resolved in protein coordinates. In addition, \dimval{VARIANT}
itself has two components: a position and a discrete nucleotide change.

We therefore record \dimval{VARIANT\_IN\_TRANSCRIPT} as a curator-surfaced
candidate value for Resolution. It is separate from the fifteen
pilot-surfaced extensions listed in \cref{tab:candidates} and awaits
independent use in a corpus paper under the promotion rule.

\section{Availability of Code, Data, and Artefacts}
\label{sec:availability}

All artefacts referenced in this paper are open-source and
available in the project repository at
\url{https://github.com/ForomePlatform/genetic-evidence-model}.
This includes the SHACL schema, the human-readable dimension
reference (\texttt{dimensions.md}), per-paper annotations for all
six pilot publications, case reports, the companion
\texttt{schema/EXTENSIONS.md}, the SHACL walkthrough
(\texttt{schema/examples.md}), the annotation and review protocol
suite (\texttt{protocols/}) and the two executable annotation and
review skills (\texttt{skills/}) described in Supplementary Note~SN3,
and the extraction tooling used to
convert PDF highlights and callouts into annotation scaffolds. The
namespace \url{https://w3id.org/genetic-evidence-model/} resolves
to the same repository.

\section*{Acknowledgements}

The idea of a semantic model for the analysis of genetic literature
grew out of 2018 conversations with Prof.\ Shamil Sunyaev; without
his early support this work could not have been completed. The
authors are also grateful to colleagues at the Brigham Genomics
Medicine (BGM) program, where much of the thinking about
clinical-genomics evidence curation took shape, and to
Vladimir Seplyarskiy for his meticulous curation of the four
manually annotated papers in the pilot corpus. This work received
no external funding.
